\documentclass[11pt]{article}
\usepackage[a4paper,margin=28mm]{geometry}
\usepackage{amsmath,amssymb,amsthm,mathtools}
\usepackage{fontspec,microtype,booktabs}
\usepackage[backend=biber,style=numeric-comp,sorting=nyt,
  giveninits=true,maxbibnames=99,doi=true,isbn=false,url=true,eprint=true]{biblatex}
\addbibresource{references.bib}
\usepackage[colorlinks=true,linkcolor=blue!55!black,citecolor=blue!55!black,
urlcolor=blue!55!black]{hyperref}
\usepackage{xcolor}
\numberwithin{equation}{section}
\newtheorem{theorem}{Theorem}[section]
\newtheorem{lemma}[theorem]{Lemma}
\newtheorem{proposition}[theorem]{Proposition}
\newtheorem{corollary}[theorem]{Corollary}
\theoremstyle{remark}
\newtheorem{remark}[theorem]{Remark}
\newcommand{\KR}{\mathsf{KR}}
\newcommand{\ord}{\operatorname{ord}}
\title{The three Kanade--Russell identities modulo nine}
\author{Yuma Mizuno}
\date{}
\begin{document}
\maketitle
\begin{abstract}
We prove the three symmetric Kanade--Russell identities modulo nine.
\end{abstract}

\section{Introduction}

Write
\[
 (a;q)_n=\prod_{j=0}^{n-1}(1-aq^j),\qquad
 (a_1,\ldots,a_s;q)_\infty=\prod_{i=1}^s\prod_{j\geq0}(1-a_iq^j).
\]
Define
\begin{equation}\label{eq:source-definition}
 \mathsf F(x,y)=\sum_{m,n\geq0}
 \frac{q^{m^2+3mn+3n^2}x^my^n}{(q;q)_m(q^3;q^3)_n},\qquad
 (\mathsf A,\mathsf B,\mathsf C)
 =\bigl(\mathsf F(1,1),\mathsf F(q,q^3),\mathsf F(q^2,q^3)\bigr).
\end{equation}
Define also
\begin{align}
 \KR_1&=\frac1{(q,q^3,q^6,q^8;q^9)_\infty},\notag\\
 \KR_2&=\frac1{(q^2,q^3,q^6,q^7;q^9)_\infty},\label{eq:K-products}\\
 \KR_3&=\frac1{(q^3,q^4,q^5,q^6;q^9)_\infty}.\notag
\end{align}

\begin{theorem}\label{thm:KR-evaluations}
As formal power series in \(q\), and as holomorphic identities on
\(|q|<1\),
\begin{equation}\label{eq:KR-evaluations}
 \mathsf A=\KR_1,\qquad \mathsf B=\KR_2,\qquad \mathsf C=\KR_3.
\end{equation}
\end{theorem}

Kanade and Russell used symbolic computation to discover new
Rogers--Ramanujan type partition identities and proposed conjectures
modulo nine and twelve in 2014 \cite{KanadeRussell2015}.
Subsequent work established several of their conjectures; in particular,
Bringmann, Jennings-Shaffer, and Mahlburg proved seven identities modulo
twelve \cite{BringmannJenningsShafferMahlburg2020}.
The three partition conjectures modulo nine corresponding to
\eqref{eq:KR-evaluations}, however,
remained unproved for more than a decade, despite the close resemblance
of their partition conditions and product forms to those of the classical
Rogers--Ramanujan identities.
For work on these conjectures, see, for example,
\cite{KanadeRussell2019,Kursungoz2019,ChernLi2020,Chern2020,
UncuZudilin2021,Konenkov2024,Tsuchioka2022,Mizuno2025,WangZhang2025}.
In this paper, we prove \eqref{eq:KR-evaluations}.

Kanade and Russell originally formulated these conjectures as equalities
between numbers of partitions \cite{KanadeRussell2015}. For \(i=1,2,3\),
let \(\mathcal D_i\) consist of partitions
\(\lambda_1\leq\cdots\leq\lambda_\ell\) whose parts are at least
\(i\), such that
\(\lambda_{j+2}-\lambda_j\geq3\), and such that
\(\lambda_j+\lambda_{j+1}\equiv0\pmod3\) whenever
\(\lambda_{j+1}-\lambda_j\leq1\).
Kanade and Russell conjectured that the
number of partitions of each integer in \(\mathcal D_i\) equals the number
with parts congruent to \(2^{i-1},3,6,9-2^{i-1}\pmod9\).
Kur\c{s}ung\"oz later expressed the generating functions for
\(\mathcal D_1,\mathcal D_2,\mathcal D_3\) as the double sums
\(\mathsf A,\mathsf B,\mathsf C\), respectively
\cite[Theorems 7--9]{Kursungoz2019}.
Thus these partition conjectures are equivalent to
\eqref{eq:KR-evaluations}.

Together with the result of Wang and Zhang
\cite[Theorem 1.1]{WangZhang2025}, Theorem~\ref{thm:KR-evaluations}
also implies that the suitably normalized triple
\((\mathsf A,\mathsf B,\mathsf C)\) is a weight-zero vector-valued
modular function on \(\Gamma_0(3)\).

The proof shows that the triple of sums \((\mathsf A,\mathsf B,\mathsf C)\)
and the triple of products \((\KR_1,\KR_2,\KR_3)\) give the same value
when substituted into the cubic polynomial
\[
 \mathcal N(X,Y,Z)=X^3+qY^3-q^2Z^3+3qXYZ.
\]
Together with the coefficientwise lower bounds
\(\mathsf A\geq\KR_1\), \(\mathsf B\geq\KR_2\), and
\(\mathsf C\geq\KR_3\), this equality forces the two triples to coincide.
Sections~\ref{app:addition}--\ref{app:Morita} establish the
identities needed to evaluate this polynomial at \((\mathsf A,\mathsf B,\mathsf C)\).
Section~\ref{sec:product} evaluates the cubic at the products using
the Jacobi addition formula and classical identities for theta and Lambert series. Section~\ref{sec:positive-completion}
combines the resulting equal cubic values with Tsuchioka's proved
coefficientwise inequalities. The first coefficient at which the
triples could differ gives a contradiction.

\paragraph*{Acknowledgments.}
The author was supported by the Irish Research Council Advanced Laureate Award
IRCLA/2023/1934 held by Robert Osburn.
The author used GPT-6 Astra in searching for the proof and preparing the manuscript.
The process is described in the final section, \nameref{sec:ai-discovery}.

\section{Bilinear identities}
\label{app:addition}

Retain the definition of \(\mathsf F\) in
\eqref{eq:source-definition}, and define
\[
 \mathsf G(x,y)=\sum_{m,n\geq0}
 \frac{q^{m^2-3mn+3n^2}x^my^n}{(q;q)_m(q^3;q^3)_n}.
\]

\begin{proposition}[Bilinear identities]\label{prop:app-addition}
The following identities hold.
\begin{align}
 \mathsf G(x,1)
 & =\mathsf F(x,1)\mathsf F(x,x^3)
       +qx^2\mathsf F(qx,q^3)\mathsf F(q^2x,q^3x^3),
       \label{eq:app-addition-0}\\
 \mathsf G(qx,1)
 & =\mathsf F(x,1)\mathsf F(qx,q^3x^3)
       -qx\mathsf F(q^2x,q^3)\mathsf F(q^2x,q^3x^3),
       \label{eq:app-addition-plus}\\
 \mathsf G(q^{-1}x,q^3)
 & =x\mathsf F(qx,q^3)\mathsf F(qx,q^3x^3)
       +\mathsf F(q^2x,q^3)\mathsf F(x,x^3).
       \label{eq:app-addition-minus}
\end{align}
\end{proposition}

\begin{proof}
For each integer \(j\), define
\[
 R_j=\sum_{n\geq0}\frac{q^{3n^2+3jn}}{(q^3;q^3)_n}.
\]
Termwise subtraction and a shift of \(n\) give
\begin{equation}\label{eq:app-R-contiguity}
 R_j-R_{j+1}=q^{3j+3}R_{j+2}.
\end{equation}
If \(g(x)=\mathsf G(x,1)=\sum_{j\geq0}g_jx^j\), then
\(g_j=q^{j^2}R_{-j}/(q;q)_j\).  Equation
\eqref{eq:app-R-contiguity} implies, for \(j\geq2\),
\[
 (1-q^j)(1-q^{j-1})g_j
 =q^{2j-1}(1-q^{j-1})g_{j-1}+q^{j-1}g_{j-2}.
\]
Equivalently,
\begin{equation}\label{eq:app-scalar-equation}
 qg(x)-(1+q+q^2x^2)g(qx)
       +(1-q^2x)g(q^2x)+q^2xg(q^3x)=0.
\end{equation}
For \(j\geq2\), the factor multiplying the coefficient of \(x^j\)
in this equation is \(q(1-q^j)(1-q^{j-1})\ne0\).
Thus a solution is determined by its constant and linear coefficients.  For \(g\),
these are
\begin{equation}\label{eq:app-g-initial}
 g(0)=R_0,\qquad [x]g(x)=\frac{q(R_0+R_1)}{1-q}.
\end{equation}

The series \(\mathsf F\) satisfies the shift relations
\begin{equation}\label{eq:app-F-contiguities}
\begin{aligned}
 \mathsf F(x,y)-\mathsf F(qx,y)
       &=qx\mathsf F(q^2x,q^3y),\\
 \mathsf F(x,y)-\mathsf F(x,q^3y)
       &=q^3y\mathsf F(q^3x,q^6y).
\end{aligned}
\end{equation}
Define
\[
 a(x)=\mathsf F(x,1),\quad b(x)=\mathsf F(qx,q^3),\quad
 \mathbf f(x)=(a(x),b(x),b(qx))^{\mathsf T}.
\]
Using \eqref{eq:app-F-contiguities} gives
\begin{equation}\label{eq:app-L-system}
 \mathbf f(qx)=L(x)\mathbf f(x),\qquad
 L(x)=\begin{pmatrix}1&0&-qx\\0&0&1\\-x&x&1+qx^2\end{pmatrix}.
\end{equation}
For example, its last row is
\(b(q^2x)=xb(x)+b(qx)-xa(qx)\).

Define \(p(x)=\mathsf F(qx,x^3)\).  It satisfies
\begin{equation}\label{eq:app-p-recurrence}
 p(x)=(1-qx)p(qx)+qx(1+q+q^2x^2)p(q^2x)+q^4x^2p(q^3x).
\end{equation}
Here is a direct verification.  Write \(p(x)=\sum c_jx^j\), with
\(c_j=0\) for \(j<0\).  The coefficient identity is
\[
 (1-q^j)c_j
 =\bigl(q^{2j-1}(1+q)-q^j\bigr)c_{j-1}
       +q^{3j-2}c_{j-2}+q^{2j-3}c_{j-3}.
\]
For each summand with \(j=m+3n\), write \(X=q^m\), \(Y=q^{3n}\).
After reindexing the three coefficients on the right by
\(m-1,m-2,n-1\), respectively, the required equality follows from the identity
\[
 1-XY+\frac{1-X}{X}-\frac{(1+q)(1-X)Y}{q}
 -\frac{(1-X/q)(1-X)Y}{X}-\frac{1-Y}{X}=0.
\]

The column vector
\(\mathbf p(x)=(p(x),p(qx),p(q^2x))^{\mathsf T}\) therefore satisfies
\begin{equation}\label{eq:app-T-system}
 \mathbf p(x)=T(x)\mathbf p(qx),\qquad
 T(x)=\begin{pmatrix}
 1-qx&qx(1+q+q^2x^2)&q^4x^2\\1&0&0\\0&1&0
 \end{pmatrix}.
\end{equation}
The shift relations also give
\begin{equation}\label{eq:app-p-bridges}
\begin{aligned}
 \mathsf F(x,x^3)&=p(x)+qx p(qx),\\
 \mathsf F(qx,q^3x^3)&=p(qx)+q^2x p(q^2x),\\
 \mathsf F(q^2x,q^3x^3)&=p(qx).
\end{aligned}
\end{equation}

Let \(h(x)\) be the right side of \eqref{eq:app-addition-0}.
Then \(h(x)=\mathbf f(x)^{\mathsf T}K_0(x)\mathbf p(x)\), where
\[
 K_0=\begin{pmatrix}1&qx&0\\0&qx^2&0\\0&0&0\end{pmatrix}.
\]
Define two further polynomial matrices by
\[
 K_1=\begin{pmatrix}0&1&q^2x\\0&0&0\\0&-qx&0\end{pmatrix},\qquad
 K_2=\begin{pmatrix}1&qx-1&-q^2x\\
                         0&0&-q^3x^3\\-qx&qx-q^2x^2&0\end{pmatrix}.
\]
Direct multiplication gives the three polynomial identities
\begin{equation}\label{eq:app-bilinear-certificate}
\begin{aligned}
 K_1(x)T(x)&=L(x)^{\mathsf T}K_0(qx),\\
 K_2(x)T(x)&=qxL(x)^{\mathsf T}K_1(qx),\\
 \bigl(q^2xK_0-(qx)(1+q+q^2x^2)K_1
       &+(1-q^2x)K_2\bigr)T(x)
       +qxL(x)^{\mathsf T}K_2(qx)=0.
\end{aligned}
\end{equation}
Contracting the first two identities with
\(\mathbf f(x)^{\mathsf T}\) and \(\mathbf p(qx)\), and using
\eqref{eq:app-L-system} and \eqref{eq:app-T-system}, gives
\[
 h(qx)=\mathbf f(x)^{\mathsf T}K_1(x)\mathbf p(x),\qquad
 qx h(q^2x)=\mathbf f(x)^{\mathsf T}K_2(x)\mathbf p(x).
\]
The same contraction of the last identity gives \(qx\) times the
left side of \eqref{eq:app-scalar-equation} with \(g=h\).
Cancelling \(qx\) shows that \(h\) satisfies that equation.

The defining sums give
\[
 a(0)=R_0,\quad [x]a=\frac{qR_1}{1-q},\qquad
 \mathsf F(x,x^3)=1+\frac{q}{1-q}x+O(x^2).
\]
The other summand in \(h\) starts at \(x^2\).  Hence \(h\) has
the two initial coefficients in \eqref{eq:app-g-initial}.
The initial two coefficients and the recurrence agree, so \(h=g\).

The matrix \(K_1\) and \eqref{eq:app-p-bridges} now give
\eqref{eq:app-addition-plus}. For the last identity, termwise
subtraction gives \(g(x)-g(qx)=qx\mathsf G(q^2x,q^{-3})\).
The relation \(R_{-j-1}=R_{-j}+q^{-3j}R_{1-j}\), from
\eqref{eq:app-R-contiguity}, then gives
\[
 qx\mathsf G(q^{-1}x,q^3)=g(x)-g(qx)-qxg(q^2x).
\]
Finally, direct multiplication gives
\[
 K_0-K_1-K_2
       =qx\begin{pmatrix}0&0&0\\0&x&q^2x^2\\1&qx&0\end{pmatrix}.
\]
Contracting with \(\mathbf f(x)^{\mathsf T}\) and \(\mathbf p(x)\),
using \eqref{eq:app-p-bridges}, and cancelling \(qx\) proves
\eqref{eq:app-addition-minus}.
\end{proof}

At \(x=1\), Proposition~\ref{prop:app-addition} yields
\begin{equation}\label{eq:app-quadratic-bridge}
 \mathsf G(1,1)=\mathsf A^2+q\mathsf B\mathsf C,\qquad
 \mathsf G(q^{-1},q^3)=\mathsf B^2+\mathsf A\mathsf C,\qquad
 \mathsf G(q,1)=\mathsf A\mathsf B-q\mathsf C^2.
\end{equation}
\begin{remark}
Expressions for the three sums on the left in terms of infinite products were conjectured
by Wang and Wang~\cite[Conjecture 3.7]{WangWang2026}.
The identities \eqref{eq:app-quadratic-bridge}, together with Theorem~\ref{thm:KR-evaluations}, prove their conjecture.
To match Wang and Wang's third formula, we rewrite its right-hand side
\(\KR_1^2\KR_3/\KR_2-2q\KR_3^2\) as \(\KR_1\KR_2-q\KR_3^2\)
using the classical theta identity
\[
 \KR_1\KR_2^2-\KR_1^2\KR_3+q\KR_2\KR_3^2=0.
\]
\end{remark}

\section{Casoratian}
\label{app:casoratian}

Define the \(q\)-Airy function with base \(p\) by
\[
 \operatorname{Ai}_p(X)=\sum_{n\geq0}
       \frac{p^{n(n-1)/2}X^n}{(p^2;p^2)_n}.
\]
Its elementary difference equation is
\begin{equation}\label{eq:app-Ai-contiguity}
 \operatorname{Ai}_p(X)-\operatorname{Ai}_p(p^2X)
       =X\operatorname{Ai}_p(pX).
\end{equation}
Consequently the function
\[
 \omega(X)=\operatorname{Ai}_p(X)\operatorname{Ai}_p(-pX)
       +\operatorname{Ai}_p(-X)\operatorname{Ai}_p(pX)
\]
satisfies \(\omega(pX)=\omega(X)\). Comparing coefficients shows that \(\omega\) is constant, hence
\begin{equation}\label{eq:app-Ai-Wronskian}
 \omega(X)=\omega(0)=2.
\end{equation}

Let \(t\) satisfy \(t^4=q\), and define
\[
 \mathcal U(t;x)=\sum_{m,n\geq0}
 \frac{(-1)^m t^{m^2+6mn+3n^2+6n}x^m}
      {(t^4;t^4)_m(t^{12};t^{12})_n},
\]
\[
 \mathcal A(x)=\frac{\mathcal U(t;x)+\mathcal U(-t;x)}2,
 \qquad
 \mathcal B(x)=\frac{\mathcal U(t;x)-\mathcal U(-t;x)}{2t},
 \qquad \mathbf r(x)=(\mathcal A(x),\mathcal B(x)).
\]
Every exponent of \(\mathcal U\) is either \(0\) or \(1\pmod4\).
Thus \(\mathcal A,\mathcal B\) are power series in \(q=t^4\) and \(x\).
Define the two rows
\begin{equation}\label{eq:app-parity-sampling}
 (A_1,B_1)=\mathbf r(1),\qquad
 (A_2,B_2)=\mathbf r(q).
\end{equation}

Let \(J_m=\operatorname{Ai}_{t^6}(t^{6m+9})\).  Then
\([x^m]\mathcal U(t;x)=(-1)^m t^{m^2}J_m/(q;q)_m\).
Equation \eqref{eq:app-Ai-contiguity} gives
\(J_{m-2}-J_m=t^{6m-3}J_{m-1}\), so these coefficients satisfy
the same recurrence as \(g_m\) in Section~\ref{app:addition}.
Both \(\mathcal A\) and \(\mathcal B\) therefore satisfy
\eqref{eq:app-scalar-equation}.  Moreover,
\eqref{eq:app-Ai-Wronskian} at \(p=t^6,X=t^9\) gives
\begin{equation}\label{eq:app-parity-initial-determinant}
 \det\begin{pmatrix}\mathbf r(0)\\\mathbf r'(0)\end{pmatrix}
 =-\frac{
 \operatorname{Ai}_{t^6}(t^9)\operatorname{Ai}_{t^6}(-t^{15})
 +\operatorname{Ai}_{t^6}(-t^9)\operatorname{Ai}_{t^6}(t^{15})}
 {2(1-q)}=-\frac1{1-q}.
\end{equation}

For each positive integer $j$, define
\[
 W_j(x)=\det\begin{pmatrix}\mathbf r(x)\\\mathbf r(q^jx)\end{pmatrix}.
\]
Equation \eqref{eq:app-parity-initial-determinant} gives
\(W_1(x)=x+O(x^2)\). Put \(H(x)=W_1(x)/x\).
Both components of \(\mathbf r\) satisfy \eqref{eq:app-scalar-equation}, so
\[
 q\mathbf r(x)-(1+q+q^2x^2)\mathbf r(qx)
 +(1-q^2x)\mathbf r(q^2x)+q^2x\mathbf r(q^3x)=0.
\]
Taking the left-hand side as the first row and \(\mathbf r(qx)\), \(\mathbf r(q^2x)\)
as the second row, respectively, and expanding the determinants gives
\begin{align*}
 q^2x W_2(qx)&=qW_1(x)-(1-q^2x)W_1(qx),\\
 qW_2(x)&=(1+q+q^2x^2)W_1(qx)+q^2xW_1(q^2x).
\end{align*}
Eliminating \(W_2\), substituting \(W_1(x)=xH(x)\),
and cancelling \(qx\) gives
\begin{equation}\label{eq:app-exterior-equation}
 H(x)=(1-q^2x)H(qx)
 +q^2x(1+q+q^4x^2)H(q^2x)+q^6x^2H(q^3x).
\end{equation}
For each positive degree \(j\), the new coefficient is multiplied
by \(1-q^j\ne0\).  Thus the coefficients of \(H\) are determined successively by its constant term.
Equation \eqref{eq:app-p-recurrence}, with \(x\) replaced by \(qx\),
shows that \(\mathsf F(q^2x,q^3x^3)\) is this solution.  Therefore
\begin{equation}\label{eq:app-Casoratian}
 \det\begin{pmatrix}
 \mathcal A(x)&\mathcal B(x)\\
 \mathcal A(qx)&\mathcal B(qx)
 \end{pmatrix}
 =x\mathsf F(q^2x,q^3x^3).
\end{equation}

In particular, at \(x=1\),
\begin{equation}\label{eq:source-minors}
 \det\begin{pmatrix}A_1&B_1\\A_2&B_2\end{pmatrix}=\mathsf F(q^2,q^3)=\mathsf C.
\end{equation}

\section{Connection formula}
\label{app:Morita}

Define the theta function by
\[
 \vartheta_p^{+}(z)=\sum_{k\in\mathbb Z}p^{k(k-1)/2}z^k.
\]
Define the Ramanujan function by
\[
 \mathcal R_q(Z)=\sum_{n\geq0}\frac{q^{n^2}(-Z)^n}{(q;q)_n}.
\]
The connection formula \cite[Section~1, Theorem]{Morita2011}, with the
\(q\)-Airy convention of Section~\ref{app:casoratian}, is
\begin{equation}\label{eq:app-Morita}
 \vartheta_p^{+}(X/p)\operatorname{Ai}_p(-X)
 +\vartheta_p^{+}(-X/p)\operatorname{Ai}_p(X)
 =2(p^2;p^2)_\infty\mathcal R_{p^2}(-p^3/X^2).
\end{equation}
The factor in the cited formula is
\((p,-1;p)_\infty=2(p^2;p^2)_\infty\).

The definitions of the parity rows give
\begin{equation}\label{eq:parity-recombination}
 \mathcal U(\pm t;1)=A_1\pm tB_1,\qquad
 \mathcal U(\pm t;q)=A_2\pm tB_2.
\end{equation}
To evaluate \(\mathcal U(t;q^j)\) for \(j=0,1\), first use
\(p=t^2\), \(X=t^{6n+4j+1}\).
Quasiperiodicity gives
\[
 \vartheta_{t^2}^{+}(\pm t^{6n+4j-1})
       =(\pm1)^n L_n\vartheta_{t^2}^{+}(\pm t^{4j-1}),
 \qquad L_n=t^{-9n^2+(6-12j)n}.
\]
Multiply \eqref{eq:app-Morita} by
\(t^{3n^2+6n}/((t^{12};t^{12})_n L_n)\), and sum over \(n\geq0\).
The exponent on the right is
\[
 4m^2-12mn+12n^2+(4-8j)m+12jn.
\]
At \(j=0,1\), respectively, it gives
\(\mathsf G(q,1)\) and \(\mathsf G(q^{-1},q^3)\).
The symmetry
\(\vartheta_{t^2}^{+}(\pm t^3)
 =\vartheta_{t^2}^{+}(\pm t^{-1})\) therefore gives
\begin{align}
 \vartheta_{t^2}^{+}(t^{-1})\mathcal U(t;1)
 +\vartheta_{t^2}^{+}(-t^{-1})\mathcal U(-t;1)
       &=2(q;q)_\infty\mathsf G(q,1),\notag\\
 \vartheta_{t^2}^{+}(t^{-1})\mathcal U(t;q)
 +\vartheta_{t^2}^{+}(-t^{-1})\mathcal U(-t;q)
       &=2(q;q)_\infty\mathsf G(q^{-1},q^3).
       \label{eq:app-short-Morita}
\end{align}

Next use \(p=t^6\), \(X=t^{6m+9}\).  In this case
\[
 \vartheta_{t^6}^{+}(\pm t^{6m+3})
       =(\pm1)^m t^{-3m^2}\vartheta_{t^6}^{+}(\pm t^3).
\]
Multiply by \(t^{4m^2+4jm}/(t^4;t^4)_m\), and sum over \(m\geq0\).
The exponent on the right is
\(4m^2-12mn+12n^2+4jm\).  Thus
\begin{align}
 \vartheta_{t^6}^{+}(-t^3)\mathcal U(t;1)
 +\vartheta_{t^6}^{+}(t^3)\mathcal U(-t;1)
       &=2(q^3;q^3)_\infty\mathsf G(1,1),\notag\\
 \vartheta_{t^6}^{+}(-t^3)\mathcal U(t;q)
 +\vartheta_{t^6}^{+}(t^3)\mathcal U(-t;q)
       &=2(q^3;q^3)_\infty\mathsf G(q,1).
       \label{eq:app-long-Morita}
\end{align}

Define
\[
 \theta(q)=\sum_{n\in\mathbb Z}q^{n^2},\qquad
 \varphi(q)=\sum_{n\in\mathbb Z}q^{n(n-1)}.
\]
Even and odd index separation gives
\[
 \vartheta_{t^2}^{+}(\pm t^{-1})
       =\varphi(q)\pm t^{-1}\theta(q),\qquad
 \vartheta_{t^6}^{+}(\pm t^3)
       =\theta(q^3)\pm t^3\varphi(q^3).
\]
Using \eqref{eq:parity-recombination} in
\eqref{eq:app-short-Morita}--\eqref{eq:app-long-Morita} consequently
proves
\begin{equation}\label{eq:app-source-matrix}
 \begin{pmatrix}A_1&B_1\\A_2&B_2\end{pmatrix}
 \begin{pmatrix}\varphi(q)&\theta(q^3)\\
                 \theta(q)&-q\varphi(q^3)\end{pmatrix}
 =\begin{pmatrix}
 (q;q)_\infty\mathsf G(q,1)&(q^3;q^3)_\infty\mathsf G(1,1)\\
 (q;q)_\infty\mathsf G(q^{-1},q^3)&(q^3;q^3)_\infty\mathsf G(q,1)
 \end{pmatrix}.
\end{equation}

\section{A cubic identity for the products}\label{sec:product}

Define the theta series \(\mathfrak a(q)\) and the Lambert series \(L(q)\) by
\[
 \mathfrak a(q)=\sum_{r,s\in\mathbb Z}q^{r^2+rs+s^2},\qquad
 L(q)=\sum_{n\ge1}\frac{q^n}{(1-q^n)^2}.
\]
We use the square identity \cite[Theorem 7.2]{Matsuda2020}
\begin{equation}\label{eq:formal-square}
 \mathfrak a(q)^2=1+12L(q)-36L(q^3)
\end{equation}
and the trisection identity \cite[equation (2.4), with $q$ replaced by $q^3$]{BorweinBorweinGarvan1994}
\begin{equation}\label{eq:formal-trisection}
 (\mathfrak a(q)-\mathfrak a(q^3))(q^3;q^3)_\infty
 =6q(q^9;q^9)_\infty^3.
\end{equation}

Write the Jacobi product as
\begin{equation}\label{eq:formal-jacobi}
 j(z;q)=(q,z,q/z;q)_\infty.
\end{equation}
It satisfies
\begin{equation}\label{eq:formal-jacobi-shifts}
 j(qz;q)=-z^{-1}j(z;q),\qquad
 j(z^{-1};q)=-z^{-1}j(z;q),\qquad j(q/z;q)=j(z;q).
\end{equation}

The Weierstrass addition formula gives
\begin{equation}\label{eq:formal-addition}
 \begin{aligned}
 &j(xz;\rho)j(x/z;\rho)j(y;\rho)^2
 -j(yz;\rho)j(y/z;\rho)j(x;\rho)^2\\
 &\hspace{15mm}=\frac yz j(xy;\rho)j(x/y;\rho)j(z;\rho)^2.
 \end{aligned}
\end{equation}

Define
\[
 \Sigma(u;\rho)=\sum_{n\ge0}\left(
 \frac{u\rho^n}{(1-u\rho^n)^2}
 +\frac{\rho^{n+1}/u}{(1-\rho^{n+1}/u)^2}\right).
\]
Expanding the paired factors of the product at $z=1+\varepsilon$ gives
\begin{equation}\label{eq:formal-kernel}
 j(u(1+\varepsilon);\rho)j(u/(1+\varepsilon);\rho)
 =j(u;\rho)^2\bigl(1-\Sigma(u;\rho)\varepsilon^2\bigr)
 \pmod{\varepsilon^3},
\end{equation}
since $(1-bz)(1-b/z)/(1-b)^2=1-b(z+z^{-1}-2)/(1-b)^2$ and $z+z^{-1}-2=\varepsilon^2\pmod{\varepsilon^3}$.

\begin{proposition}\label{prop:mc-product-cubic}
The three products satisfy
\begin{equation}\label{eq:mc-product-norm}
 \mathcal N(\KR_1,\KR_2,\KR_3)
 =\frac{\mathfrak a(q)}{(q;q)_\infty(q^3;q^3)_\infty}.
\end{equation}
\end{proposition}

\begin{proof}
For $1\le r<9$, write $J_r=j(q^r;q^9)$ and $\Lambda_r=\Sigma(q^r;q^9)$. Splitting the products by residue classes gives
\begin{gather*}
 J_3=(q^3;q^3)_\infty,\qquad
 J_1J_2J_4=\frac{(q;q)_\infty(q^9;q^9)_\infty^3}{(q^3;q^3)_\infty},\\
 (\KR_1,\KR_2,\KR_3)
 =\frac{(q^9;q^9)_\infty^2}{(q^3;q^3)_\infty}
 \left(\frac1{J_1},\frac1{J_2},\frac1{J_4}\right),\\
 \KR_1\KR_2\KR_3=\frac{(q^9;q^9)_\infty^3}{(q;q)_\infty(q^3;q^3)_\infty^2}.
\end{gather*}
In \eqref{eq:formal-addition}, substitute $(\rho,x,y)=(q^9,q^r,q^3)$ and $z=1+\varepsilon$. The second coefficient, using \eqref{eq:formal-kernel}, gives
\[
 J_r^2(q^3;q^3)_\infty^2(\Lambda_r-\Lambda_3)
 =-q^3(q^9;q^9)_\infty^6j(q^{r+3};q^9)j(q^{r-3};q^9).
\]
For $r=1,2,4$, reflection and the preceding product relations simplify these identities to
\begin{align*}
 \Lambda_1-\Lambda_3&=q(q;q)_\infty(q^9;q^9)_\infty^3\KR_1^3,\\
 \Lambda_2-\Lambda_3&=q^2(q;q)_\infty(q^9;q^9)_\infty^3\KR_2^3,\\
 \Lambda_4-\Lambda_3&=-q^3(q;q)_\infty(q^9;q^9)_\infty^3\KR_3^3.
\end{align*}
Their sum has left side $L(q)-4L(q^3)+3L(q^9)$. Multiplying by $(q^3;q^3)_\infty^2$ and using the product of the three $\KR_i$ gives
\begin{equation}\label{eq:formal-norm-reduction}
 \begin{aligned}
 &(q^3;q^3)_\infty^2\bigl(L(q)-4L(q^3)+3L(q^9)\bigr)\\
 &\quad=q(q;q)_\infty(q^3;q^3)_\infty^2(q^9;q^9)_\infty^3
 \mathcal N(\KR_1,\KR_2,\KR_3)
 -3q^2(q^9;q^9)_\infty^6.
 \end{aligned}
\end{equation}

It remains to evaluate the Lambert expression. Subtracting \eqref{eq:formal-square} at $q^3$ from the same identity at $q$ yields
\[
 12\bigl(L(q)-4L(q^3)+3L(q^9)\bigr)
 =(\mathfrak a(q)-\mathfrak a(q^3))(\mathfrak a(q)+\mathfrak a(q^3)).
\]
Applying \eqref{eq:formal-trisection} and writing the sum of the two theta series as $2\mathfrak a(q)-(\mathfrak a(q)-\mathfrak a(q^3))$ gives
\[
 (q^3;q^3)_\infty^2\bigl(L(q)-4L(q^3)+3L(q^9)\bigr)
 =q\mathfrak a(q)(q^3;q^3)_\infty(q^9;q^9)_\infty^3
 -3q^2(q^9;q^9)_\infty^6.
\]
Comparison with \eqref{eq:formal-norm-reduction} and cancellation proves the proposition.
\end{proof}

\section{Proof of the main theorem}
\label{sec:positive-completion}

Taking determinants in \eqref{eq:app-source-matrix} and applying
Proposition~\ref{prop:mc-product-cubic} gives the following equalities.

\begin{lemma}[Equality of the cubic norms]\label{lem:equal-norms}
As identities of formal power series,
\begin{equation}\label{eq:equal-norms}
 \mathcal N(\mathsf A,\mathsf B,\mathsf C)
 =\frac{\mathfrak a(q)}{(q;q)_\infty(q^3;q^3)_\infty}
 =\mathcal N(\KR_1,\KR_2,\KR_3).
\end{equation}
\end{lemma}

\begin{proof}
Splitting the summation lattice according to the parity of \(s\)
gives
\begin{equation}\label{eq:a2-parity}
 \mathfrak a(q)=\theta(q)\theta(q^3)
                   +q\varphi(q)\varphi(q^3).
\end{equation}
Indeed, the substitutions \((r,s)=(\ell-k,2k)\) and
\((r,s)=(\ell-k,2k-1)\) give the exponents
\(\ell^2+3k^2\) and
\(\ell(\ell-1)+3k(k-1)+1\), respectively. Taking determinants directly in
\eqref{eq:app-source-matrix}, and using \eqref{eq:source-minors}
and \eqref{eq:a2-parity}, gives
\[
 \mathsf G(q,1)^2-\mathsf G(1,1)\mathsf G(q^{-1},q^3)
 =-\frac{\mathsf C\mathfrak a(q)}{(q;q)_\infty(q^3;q^3)_\infty}.
\]
Substituting \eqref{eq:app-quadratic-bridge} yields
\[
 (\mathsf A\mathsf B-q\mathsf C^2)^2
 -(\mathsf A^2+q\mathsf B\mathsf C)(\mathsf B^2+\mathsf A\mathsf C)
 =-\frac{\mathsf C\mathfrak a(q)}{(q;q)_\infty(q^3;q^3)_\infty}.
\]
The polynomial identity
\begin{equation}\label{eq:norm-determinant-factorization}
 (XY-qZ^2)^2-(X^2+qYZ)(Y^2+XZ)
       =-Z\mathcal N(X,Y,Z)
\end{equation}
therefore proves the first equality in \eqref{eq:equal-norms}.
The cancellation of \(\mathsf C\) is valid because its constant coefficient is \(1\).

The second equality is Proposition~\ref{prop:mc-product-cubic}.
\end{proof}

One further input supplies the direction of every possible
discrepancy between the two triples.

\begin{lemma}[Coefficientwise lower bounds]\label{lem:KR-lower-bounds}
The three differences satisfy
\begin{equation}\label{eq:KR-lower-bounds}
 \mathsf A-\KR_1,\quad \mathsf B-\KR_2,\quad
 \mathsf C-\KR_3\ \in q\mathbb Z_{\geq0}[[q]].
\end{equation}
\end{lemma}

\begin{proof}
Let \(\mathcal D_i\), \(i=1,2,3\), be the partition sets in the introduction.
Kur\c{s}ung\"oz proves that the generating functions
\(\sum_{\lambda\in\mathcal D_i}q^{|\lambda|}\) are
\(\mathsf A,\mathsf B,\mathsf C\), respectively
\cite[Theorems 7--9, equations (1), (38), (85)]{Kursungoz2019}.
In his two-variable formulas, the length variable is set to \(1\).
The three exponents, in the present indices \((m,n)\), are
\[
 m^2+3mn+3n^2,\qquad
 m^2+3mn+3n^2+m+3n,\qquad
 m^2+3mn+3n^2+2m+3n.
\]
These are exactly the exponents in the definitions of \(\mathsf A,\mathsf B,\mathsf C\).

Tsuchioka's spanning theorem implies that, for every \(N\geq0\),
the number of partitions of \(N\) in \(\mathcal D_i\) is at least
the number of partitions of \(N\) into parts congruent to
\(2^{i-1},3,6,9-2^{i-1}\pmod9\)
\cite[Corollary 1.4]{Tsuchioka2022}.
The latter generating functions are \(\KR_1,\KR_2,\KR_3\).
All six series have constant coefficient \(1\), proving
\eqref{eq:KR-lower-bounds}.
\end{proof}

\begin{lemma}[Rigidity under coefficientwise lower bounds]
\label{lem:positive-norm-rigidity}
Let \(X,Y,Z,x,y,z\in1+q\mathbb Z[[q]]\).  If
\(X-x,Y-y,Z-z\) have nonnegative coefficients and
\(\mathcal N(X,Y,Z)=\mathcal N(x,y,z)\), then
\((X,Y,Z)=(x,y,z)\).
\end{lemma}

\begin{proof}
Write \(d_1=X-x\), \(d_2=Y-y\), and \(d_3=Z-z\).
Telescoping the cubic and mixed terms gives the exact identity
\begin{equation}\label{eq:positive-norm-factorization}
\begin{aligned}
 \mathcal N(X,Y,Z)-\mathcal N(x,y,z)
 ={}&d_1\bigl(X^2+Xx+x^2+3qYZ\bigr)\\
 &+qd_2\bigl(Y^2+Yy+y^2+3xZ\bigr)\\
 &+qd_3\bigl(3xy-q(Z^2+Zz+z^2)\bigr).
\end{aligned}
\end{equation}
The three expressions in parentheses have constant coefficients
\(3,6,3\), respectively.  Suppose that at least one \(d_i\) is
nonzero, and define
\[
 \nu=\min\{\ord_qd_1,\ 1+\ord_qd_2,\ 1+\ord_qd_3\},
 \qquad \ord_q0=+\infty.
\]
Each summand of order \(\nu\) on the right of
\eqref{eq:positive-norm-factorization} has a positive coefficient
at that order.  Every other summand has greater order.  Thus the
coefficient of \(q^\nu\) is positive, contradicting the equality of
the norms.  All three differences vanish.
\end{proof}

\begin{proof}[Proof of Theorem~\ref{thm:KR-evaluations}]
Apply Lemma~\ref{lem:positive-norm-rigidity} to the equal norms in
Lemma~\ref{lem:equal-norms} and the lower bounds in
Lemma~\ref{lem:KR-lower-bounds}.
\end{proof}

\section*{Use of AI}
\label{sec:ai-discovery}

I would first like to describe how I came to ask AI to solve this conjecture.
I learned about the conjecture while I was a student in the Department of Mathematical and Computing Science
at Tokyo Institute of Technology (now Institute of Science Tokyo), and I remember hearing about it from Shunsuke Tsuchioka.
I do not remember exactly when I first heard of it, but hearing about it in Tsuchioka's lectures
(on Rogers--Ramanujan identities and related topics), which I attended via Zoom in 2020, left a lasting impression on me.

When I saw the conjecture, I noticed that the $q$-series on the sum side could be realized as
partition $q$-series of Kato and Terashima~\cite{KatoTerashima2015} in the context of cluster algebras,
which I was studying at the time (Terashima was then my supervisor).
This realization could be used to derive related dilogarithm identities, but it did not provide a way to prove
the original $q$-series identities. Ever since then, the conjecture had remained on my mind.
Its connection with cluster algebras was also one of the main motivations for my later paper~\cite{Mizuno2025}.
In that paper, I stated the modular transformation formulas for the products without proof.
I had in fact obtained a proof of the transformation formulas by expressing the products as double theta series,
but had not written it up because it lay outside the main subject of the paper and I thought it might be useful
for the Kanade--Russell conjecture itself.
Wang and Zhang~\cite{WangZhang2025} later proved these transformation formulas with a shorter proof than mine,
so in the end, I think it was good that I had left my proof out.

Because this history had kept the problem on my mind, I tried asking each new AI model whether it could solve the conjecture.
At first, I gave the models my own previous attempts (mostly recorded in SageMath files) and asked whether they could solve it.
Eventually, the files generated during their attempts came to outnumber my own.
The first model to make meaningful progress was GPT-5.6 Sol, which discovered on July 22, 2026 that the series
on the sum side could be expressed as a Casoratian, as in \eqref{eq:app-Casoratian}.
This did not yet lead to a complete proof.
When GPT-6 Astra became available in September 2026, I asked, as usual, whether they could solve the conjecture.
This time, the attempt led to a complete proof.
The key discovery was the equality of cubic norms in Lemma~\ref{lem:equal-norms}.
A clue came from Wang and Wang's conjectural expressions in terms of infinite products~\cite[Conjecture 3.7]{WangWang2026}:
after simplification, their right-hand sides are the three quadratic expressions
in \eqref{eq:app-quadratic-bridge}, with the products \(\KR_1,\KR_2,\KR_3\) in place of the sums.
This suggested the identities \eqref{eq:app-quadratic-bridge}.
GPT-6 Astra proved them via Proposition~\ref{prop:app-addition} and combined them with the connection formula
and the Casoratian identity; the determinant factorization \eqref{eq:norm-determinant-factorization}
then produced the cubic \(\mathcal N\).
This cubic already arises naturally from the modular transformation formulas for the products
that I had stated in the paper mentioned above, but I had not thought to consider it.

The reason this cubic arises naturally can be explained as follows. Put $q=e^{2\pi i\tau}$ with $\operatorname{Im}\tau>0$,
$\zeta=e^{2\pi i/9}$, and $\omega=\zeta^3$. The column vector
\[
 f(\tau)=(q^{-1/18}\KR_1,q^{5/18}\KR_2,q^{11/18}\KR_3)^{\mathsf T}
\]
satisfies $f(-1/(3\tau))=Sf(\tau)$~\cite{Mizuno2025,WangZhang2025}, where
\[
 S=\begin{pmatrix}
 \alpha_1&\alpha_2&\alpha_4\\
 \alpha_2&-\alpha_4&-\alpha_1\\
 \alpha_4&-\alpha_1&\alpha_2
 \end{pmatrix},\qquad
 \alpha_j=\frac{1}{2\sqrt3\sin(j\pi/9)}.
\]
For the linear forms $L_j:\mathbb C^3\to\mathbb C$,
$L_j(X,Y,Z)=X+\omega^jY-\omega^{2j}Z$ $(j=0,1,2)$, we have
\[
 \begin{gathered}
 L_0L_1L_2=X^3+Y^3-Z^3+3XYZ,\\
 L_0\circ S=L_0,\qquad
 L_1\circ S=\zeta L_2,\qquad
 L_2\circ S=\zeta^{-1}L_1.
 \end{gathered}
\]
Thus their product is invariant under $S$. Evaluating it at $f(\tau)$ gives
\[
 L_0(f(\tau))L_1(f(\tau))L_2(f(\tau))
 =q^{-1/6}\mathcal N(\KR_1,\KR_2,\KR_3).
\]
This recovers the cubic $\mathcal N$.
These transformation laws, combined with a comparison of the expansions at the cusps,
also give a modular proof of Proposition~\ref{prop:mc-product-cubic}.

GPT-6 Astra found that the same value could also be obtained from the sums, establishing a clear relation between the sum and product sides.
They then combined this equality with Tsuchioka's coefficientwise inequalities
(Lemma~\ref{lem:KR-lower-bounds}) as the final piece of the proof, thereby proving that the two sides coincide.

\printbibliography
\bigskip

\noindent
\textsc{Yuma Mizuno, School of Mathematical Sciences, University College Cork, Western Road,
Cork, Ireland.}\par
\noindent
\textit{Email address}: \texttt{YMizuno@ucc.ie}

\end{document}
